\vspace{-0.1em}
\section{Overview of Proposed Approach}
In this paper we propose an intent-aware deep reinforcement learning (DRL)-based technique: \name{} that can make the optimal choice of samples to be used at each step of the analysis flow to speed interactions but prevents intent-divergence. 
%For example, in Figure~\ref{fig:intent_divergence}(a.2) an intelligent sample selection agent at the back-end would have decided to pick a \textit{diversified sample} (that would preserve large delay values for \texttt{JUN}) or a \textit{stratified sample} (that would include higher proportion of rows for \texttt{JUN}, if there were less number of entries in the original data). 

\begin{comment}
\subsection{Problem Formulation}
\label{sec:problem}
%
%\sm{Please review. Does this problem formulation make sense or does it becomes confusing or conflicting}

%
\tong{Is $D$ the variable $A$ in Section 4.1?}
\sm{Yes}
Let $D_0$ be the full data target for EDA.
Let $D'$ = $\{D_i\}_{i=1}^n$ be $n$ subsamples created from $D$ and $|D_i| < |D_0|$ for $i \in [1, n]$ where $|.|$ denotes size (\# of rows). 
Let an $l$ length EDA session $e=\langle \{q_j(D_i), v_j(D_i)\}_{j=1}^l | i \in [1,n]\rangle$ be comprised of an alternating sequence of queries $q_j(.)$ and corresponding results/visual output $v_j(.)$ performed on data $D_i$. 
%

Let $I$ be the set of implicit intents embedded in the distribution of all $e$, the EDA sessions.
%\sg{What is e here?}\sm{e is an EDA sequence/session}
%
Let $I^{D_0}$ be original distribution of the intents $P(I|D_0)$, i.e. when all EDA was done using full data $D_0$.
Let $I^{D'}$ be the distribution $P(I|D')$, i.e., EDA is approximated using only subsamples.
Our goal is to find optimal subsample $D_i, i \in [1,n]$ for each step of the EDA $\{q_j(.), v_j(.)\}$ such that $Div(I^{D_0}||I^{D'})$ is minimized with $I$ as the support.
%

We define $Div(I^{D_0}||I^{D'})$ as the \textbf{intent-divergence} between the EDA performed using full data and approximate EDA using only subsampled data.  \tong{if we formally define this metric, is this metric reflected in or related to some evaluation metrics in Section 5.3? If not, might it be better not to formally define this metric?}

%\sm{We should clarify intent-divergence here.}
\end{comment}


\subsection{Design of \name{}}
%\vspace{-0.2em}
We show the high-level system architecture of \name{} in Figure~\ref{fig:intent_divergence}(b).
\name{} first takes in a history of EDA sequences performed by human data analysts and uses Topic Modeling technique to identify a set of implicate intents in those sequences. 
Then the RL-agent of \name{} takes into account:
(1) these clusters of these latent intents,
(2) the history of the ongoing EDA exploration session including the query used so far and the corresponding display-outputs (graphs and dataframes),
(3) the next query the user is intending to run,
and (4) the set of available samples created with different sampling strategies along with the size of each. 
The RL-agent, which is parameterized by a deep neural network~\cite{lillicrap2015, mnih2015human}, is trained \textit{offline} to choose the optimal sampling strategy as the best \textit{action} for different context of the analyses and intent. We considered alternate RL algorithms (DQN [\cite{mnih2013playing}] and REINFORCE [\cite{sutton2018reinforcement}]) but found the convergence and stabilization of A2C training (Lillicrap et. al), which is a hybrid between value and policy based approach, was better and reduces variance.
The best action corresponding to each step, i.e. for each query in the EDA session attempts to minimize the divergence of intents due to approximation error caused by different samples, while optimizing for lower latency. 

%\tong{}

%\noindent \textbf{Why DRL?}
%In our system, we choose the optimal sampling strategy by DRL. In RL, at a high-level, an agent interacts with a system and tries to learn an optimized policy.
%In DRL, neural networks are used as agents to handle large state and action spaces~\cite{lillicrap2015, mnih2015human}.
For our problem, lack of explicit intents in the EDA sessions, a moderately large action space of possible sampling strategies combined with the choice of parameters for each of those strategies that would control the statistical information preserved in the samples and associated trade-off for latency based on the size of the samples ($\#$ of rows), makes this a complex optimization problem. 
Appendix\textbf{~\ref{app:discussion}} further discusses why RL better choice instead of greedy solutions.

\subsection{Intent Identification}
\label{sec:intent_identification}



The problem of intent aware sample selection requires us to first solve the problem of identifying the intent of the user given a query sequence $Q_j$. 
%We use the set of EDA sequences generated by a simulator~\cite{atena2020sigmod} to model different intents.
%
% \tong{we may cite some intent-ware papers without using RL.}
%
%In this section, we draw an analogy between a collection of natural language documents and a sequence of queries. A word, which is a fundamental unit of a document corresponds to a query in our case. 
Similar to analyzing natural language texts to understand the topics, we can analyze query sequences to discover the user intents. Inspired by the previous work which uses topic modeling techniques to model goals
that users carry out by executing a sequence of commands \cite{aggarwal2020goal}, we leverage topic modeling to discover the user intent, which is further used to derive the reward signal as detailed later in Section \ref{sec:reward}. 

%\sm{“They however do not properly motivate Topic Modeling as a means to identify the intent of a user; Topic Modeling appears to be an arbitrary choice and is disconnected from the actual intent of the user. Topic Modelling is a very limited modelling technique and involves many delicate choices to get right even when it is appropriate. One obvious problem in the setting at hand is how many topics to choose to match the wide (potentially infinite) range of potential user intents during their exploration.”}
%\sg{https://dl.acm.org/doi/10.1145/3383313.3412255}


Specifically, we use a bi-term model (BTM) \cite{btm2013www} for the identification of the intent. Each topic output from the BTM model corresponds to an intent. For a query sequence $Q_j$, the output from the BTM model, \(\phi\) is the probability distribution over all the topics.

{\small
\begin{equation}
    \phi(Q_j) = \{P(t=t_i | Q=Q_j)\}_{i=1}^kK
\end{equation}
}
Here $K$ is the number of topics, and each $t_i$ corresponds to a topic. To get the corresponding intent, we take the $argmax$ over the probability distribution. The obtained intent, $I_{Q_j}$ is considered to be the user intent by \name{}.

{\small
\begin{equation}
    I_{Q_j} = argmax(\{P(t=t_i | Q=Q_j)\}_{i=1}^K)
\end{equation}
}
To decide the optimum number of topics for the BTM model, we use a coherence evaluation metric, the UCI measure~\cite{newman2010evaluating}, defined as

{\small
\begin{equation}
    score_{UCI}(q_i, q_j) = \log{\frac{p(q_i, q_j)}{p(q_i)p(q_j)}},
\end{equation}
}
where $p(q_i)$ represents the probability of seeing a query $q_i$ in a query sequence and $p(q_i, q_j)$ is the probability of observing both $q_i$ and $q_j$ co-occurring in a query sequence computed as follows,$p(q_i) = \frac{M(q_i)}{M} \text{ and } p(q_i, q_j) = \frac{M(q_i, q_j)}{M}$, where $M(q_i)$ is the count of query sequences containing the query $q_i$, $M(q_i, q_j)$ is the count of sequences containing both $q_i, q_j$ and $M$ is the total number of query sequences.
The $UCI$ score for an intent $I$, is calculated as, $mean\{score_{UCI}(q_i,q_j | q_i, q_j \in I, i \neq j)\}$ and the overall UCI score the average across all intents. Higher UCI score is indicative of better grouping because one would expect two queries belonging to the same intent to show up together regularly in the query sequences. We compute the overall UCI score for all $K \in [2,15]$ and choose the best K.

% \sm{We need the plot for housing as well.}

% \sm{Please add how to you choose he optimum cluster.}
% \sm{Please explain these and the plot almost similarly as the Goal-driven paper. }

\section{RL Formulation}
%\vspace{-0.2em}

%\subsection{RL Formulation and Rewards}
We use RL for intent-aware sample selection the agent learns the optimal policy modeled as a conditional distribution $\pi(a|s)$, specifying the probability of choosing action $a \in A$ when in state $s \in S$. 
If the agent chooses an action $a \in A$ at state $s \in S$, it receives a reward $r(s,a)$.

\subsection{Action Space}
\vspace{-0.2em}
A set of $n$ samples are pre-created using different sampling strategies (and of different sizes). 
These samples are the set of $n$ actions: $A = \{a_1, a_2, .., a_i\}_{i=1}^n$ for the RL-agent.

\subsection{State Space}
%\vspace{-0.2em}
As shown in Figure~\ref{fig:intent_divergence} (b), we use 4 components while deciding which action to choose. These components form the state space of the RL agent. The agent at each timestep $t$ for a query $q_t$ has access to the ongoing EDA state comprising of the previous $k$ queries $\{q_{t-1}, q_{t-2}, ..., q_{t-k}\}$ and their result vectors $\{v_{t-1}, v_{t-2}, ..., v_{t-k}\}$. 
%\sm{What is the notation for a query? small q or capital Q? I think we are mixing it up.}
%\sg{q stands for a single query and Q stands for a sequence of queries}
Since reducing the latency is one of the primary goals of the agent, it also takes into account the latency encountered till timestep $t$. We denote this latency factor by $C_t=\sum_{i=0}^t c_i$, where $c_i$ is the latency associated with running a query $q_i$ on a sample $a_i$. Finally, to preserve the intent information of the sequence, we also include the probability distribution over all the topics, $I_t = \phi(Q_t)$ as one of the state space components. Formally, the state space can be written as, 
{\small
\begin{equation}
\label{eq:state_space}
    s_t = \{ ((q_{t}, v_{t}),(q_{t-1}, v_{t-1}), (q_{t-2},v_{t-2})), I_t, C_t  \}
\end{equation}
}

Once the agent chooses an action $a_t \in A$, the agent moves to the next state $s_{t+1}$ and returns a result (display vector) $v_t$, as the result of the query $q_t$. Our goal here is to choose $a_t$ at each step, such that the overall latency for the sequence is minimised, simultaneously preserving the intent of the generated sequence. To do this we develop several reward functions and propose an intent-aware RL framework based on A2C methodology~\cite{mnih2016asynchronous}.
Details of our training algorithm is given in supplement (\ref{app:training}).
% \sm{Citation not coming.}

\subsection{Reward Design}
\label{sec:reward}
\vspace{-0.2em}
\name{} calculates the reward at the end of each exploratory session. For the calculation of reward, we use a set of sequences generated on the original table and we refer them as ground truth, $O_i \in O$. Each $O_i$ is a sequence of queries. Similarly, the sequence generated by the agent at the end of each episode is represented as $Q_t$.
Reward design combines following components.

\noindent\textbf{Latency Reward. }
Choosing samples of smaller size gives lower latency and hence this reward is calculated as:
%At a state $s_t$, when an agent takes an action $a_t$, the total number of rows read by the agent while taking that action is termed as $|{T_s}|$ and the total number of rows in the original dataset as $|T|$. Then the immediate reward after taking an action $a_t$ is,

{\small
\begin{equation}
\label{latency_reward}
    R_{latency} = \sum_{t=1}^{l} r(s_t, a_t),
\end{equation}
}
,where $r(s_t, a_t) = 1 - \frac{|{T_s}|}{|T|}$, where $T_s$ is the time taken to run on a sample, $T$ is the time taken to run on the original dataset, and $l$ is the length of the sequence in the session. 
This reward helps to navigate the trade-off between lower approximation error vs. lower latency. 
%that the agent does not always select the sample with the most amount of data, as that data usually leads to best approximation of data but high latency.

\noindent\textbf{Intent Reward. }
Intent reward is designed to prevent divergence of EDA sequences due to approximations (\S!\ref{sec:problem}). 
%While choosing the smallest samples using the latency reward will lead to reduced waiting time for the user, but it will often lead to diverging of intent after some point. If we only use the latency reward, the agent will often choose the smallest sample out of all the available ones. 
%We not only want the RL-agent to choose samples that will minimize time, but also to preserve the intent of the analyses. 
To prevent intent-divergence we introduce two components in the intent reward. We denote the distance function for distance between two query sequences corresponding to two EDA sessions to be $Distance(Q_1, Q_2)$. We define the reward function to be,

{\small
\begin{equation}
\label{eq_eda_sim}
    R_{dis} = 1 - Distance(Q_t, Q_{ground}).
\end{equation}
}
where $Q_{ground}$ is the closest sequence to the given sequence $Q_{t}$ in the original set $O$.
 We use \texttt{EDA-Sim} defined in \cite{atena2020sigmod} as distance metric. EDA-Sim compares the content and order of the results from the query sequence to calculate a similarity score. Based on this score we find a sequence from $O$ that is closest with current sequence.
%, so that we can match our trajectory with one of the sequences from the original sequences. 

However, it is possible that two sequences may not have the exact same ordering or content (resulting in low EDA-Sim score) but still overall have same underlying intent. 
%only the above intent reward is not sufficient for preserving the intent. There are cases where query sequences with different intents may have a low EDA-Sim value. 
Hence we also use Euclidean Distance ($EuD$) to compare intent identified from topic model distributions (\S~\ref{sec:intent_identification}) as follows:


{\small
\begin{equation}
\label{topic_reward}
    R_{topic} = 1 - EuD(\phi(Q_t), \phi(Q_{ground}))
\end{equation}
}

Full intent reward is: $R_{intent} = R_{dis} + \delta*R_{topic}$.
%($\delta$ is hyperparameter),

%{\small
%\begin{equation}
%\label{intent_reward}
%    R_{intent} = R_{dis} + \delta*R_{topic}
    %_{intent} = R_{dis} + R_{topic}
%\end{equation}
%}

%The above rewards ensure that the topic distribution matches with that of the closest sequence.
\begin{comment}
\noindent\textbf{Batch Intent Reward}
When using intent and latency rewards, the agent might always select samples corresponding to a only a few intent. To ensure that the agent also learns to preserve the distribution of different intents across various EDA sessions, we enforce another reward called batch-intent-reward.
%We want our agent to choose samples, so as to preserve the original distribution of intent over ground sequences $O$. Therefore, %
We add a component of reward which is the difference in consecutive differences between intent distributions. We term this reward as,
{\small
\begin{equation}
    R_{batch} = \bigl(D(U(Q_{t-1})) - D(U(O))\bigr) - \bigl(D(U(Q_{t})) - D(U(O))\bigr),
\end{equation}
}
where $U$ is the distribution of intent over the sequences.
\end{comment}

\noindent\textbf{Termination Reward. }
While the whole EDA session, comprising of query and corresponding visualization sequences, are important to the analyst, usually last few queries are the most important as they provide key takeaway insights. 
%
%Even after, matching intent of the user with the closest sequence, it might not be enough to preserve these insights. 
%
Therefore, the another component of our reward is based on final insight preservation. The first subcomponent $R_{match} = \{0, 1\}$, a binary 
reward if the last $k$ queries of the generated sequence $Q_t$ match with at least one of the ground truth sequence belonging to the same intent cluster.
%
However, often different query sequence lead to similar insights. Therefore, the other subcomponent  $R_{recall}$ as the top-k recall for the display vectors (results) from the generated sequence, $Q_t$ and the closest sequence identified, $Q_{ground}$ in Eqn.~\ref{eq_eda_sim}.
Full termination reward is: $R_{term} = R_{match} + R_{recall}$.

%With all the reward functions, the reward function $J_R(\pi)$ is
%
%{\small
%\begin{equation}
%\label{policy_equation}
%    J_R(\pi) = \mathbb{E}_{a_t \sim \pi_\theta}[R_t],
%\end{equation}
%}

Complete reward is $R_t = R_{latency} + \beta*R_{intent}+\gamma*R_{term}$.

% is,
%{\small
%\begin{equation}
%\label{total_reward}
%    R_t = R_{latency} + \beta*R_{intent}+\gamma*R_{term}.
%\end{equation}
%}
%To train the RL model with the given reward function we use the policy-gradient bases Advantage Actor Critic (A2C) with neural network as a function approximator for the policy and value functions. 
%For the policy network and the value network, we have different neural networks, each with 2 hidden layers of latent dimension 64.
%After learning the policy $\pi_\theta(a|s_t)$ by the above reward function, the action at each step is chosen via,
%{\small
%\begin{equation}
%\label{action_argmax}
%    a_t = argmax_{a \in A}\pi_\theta(a | s_t).
%\end{equation}
%}

\begin{comment}

\subsection{Problem Formulation}

\sm{+ mention run-time sampling has significant latency and complicated to implement engineering-wise}
Our task is to choose the best sampling strategy for an interactive exploratory data analysis workflow minimising the cost of running the queries while preserving the intent of the analyst performing this analysis. We pre-compute a set of $n$ samples, $A$ for the dataset given. This is done because online sampling techniques have significant latency and are complex to implement. The agent is given this set of sampled datasets with $n$ samples, \( A = \{a_1, a_2, .., a_i\}_{i=1}^n \), \tong{?} where each $a_i$ is a downsampled version of the original dataset. The computation cost array \( C = \{c_1, c_2, .., c_i\}_{i=1}^n \), where each $c_i$ is the downsampling percentage for $a_i$. \tong{unclear whether it is a linear function, shaddy-remove} The cost of computation is directly proportional to the number of rows in a dataset and we use this to train our RL agent.
\tong{it is unclear why D is a finite set. May need some explanations.}

The agent is provided with a sequence of queries of length $l$, \(Q = \{q_1, q_2, .., q_i\}_{i=1}^l\) to be run on a dataset, $a_0$. We denote the intent of the user interacting with the query engine to be $I_{orig}$ \tong{?}. We formulate this decision problem as a Markov Decision Process (MDP) \tong{when introducing MDP, you may describe in a way similar to Section 2.1 in https://proceedings.neurips.cc/paper/2019/file/52130c418d4f02c74f74a5bc1f8020b2-Paper.pdf}. At each step $t$ in the query sequence, the state of the agent is $s_t$, and the possible action space is $\{a_i\}_{i=1}^n $. The state $s_t$ can be represented by the vector of $k=2$ previous queries and their display vectors, the intent distribution from the intent identification part of our pipeline, and the latency cost uptil that point. Formally, it can be written as
\begin{equation}
\label{eq:state_space}
    s_t = \{ ((q_{t}, v_{t}),(q_{t-1}, v_{t-1}), (q_{t-2},v_{t-2})), I_t, C_t  \}
\end{equation}
\tong{?}
Here $v_i$ represents the display vector corresponding to the query $q_i$. Detailed information about this can be found in Section~{3.3}.

Our goal here is to choose $d_t$ at each step $s_t$, such that the overall cost $C_{total} = \sum_{i=0}^l c_i$ \tong{discount?} for the sequence is minimised, simultaneously preserving the intent of the generated sequence, $I_{gen}$ \tong{?} matches the original intent, $I_{orig}$. 
\tong{We may better define `intent' and `matches the original intent' somewhere in the paper.}
Once the agents picks an action $a_t \in A$, the agent moves to the next state $s_t+1$ and returns to user $v_t$, which is the result of the query $q_t$. To solve our MDP problem \tong{optimize XX over time}, we propose an intent-aware RL framework which uses A2C algorithm with intent based reward.
\end{comment}

\begin{comment}

\subsection{Intent Modelling}
\sm{Should go before Model Overview}


\sg{
\colorbox{yellow}{Experiment section 2: } - Include heatmaps for the generated sequences showing the distrubution of the queries in different clusters. 
\begin{itemize}
    \item Try various clustering sizes (6,8,10) and see what makes sense. Maybe include all the results here.
    \item Show the diverging intents on the original and fixed sampling approach.
\end{itemize}
}

The problem of intent aware sample selection requires us to first solve the problem of identifying the intent of the user given a query sequence $Q_j$. We use the set of EDA sequences generated by the simulator to model different intents.

%We draw an analogy between a collection of natural language documents and our set of generated query sequences. In a document, a word is a fundamental unit. 
\tong{we may cite some intent-ware papers without using RL.}
Similar to analyzing natural language documents to understand the topics, we can analyze the query data to discover the user intents\cite{goal2020recsys}
A query $q_i$ is a fundamental unit for the construction of a query sequence. We use a bi-term model (BTM) \cite{btm2013www} for the identification of the intent. Each topic output from the BTM model corresponds to an intent. For a query sequence $Q_j$, the output from the BTM model, \(\phi\) is the probability distribution over all the topics.
\begin{equation}
    \phi(Q_j) = \{P(t=t_i | Q=Q_j)\}_{i=1}^k 
\end{equation}
Here $k$ is the number of topics, and each $t_i$ corresponds to the topic. For getting the intent, we take the $argmax$ over the probability distribution
\begin{equation}
    I_{Q_j} = argmax(\{P(t=t_i | Q=Q_j)\}_{i=1}^k)
\end{equation}
This $I_{Q_j}$ is considered to be the intent for our problem.

\end{comment}

\begin{comment}
\subsection{Model Overview}

\subsubsection{State-Space Encoding}
Our model takes in a query, $q_t$ from a user at time step $t$. We support three types of EDA operations ($OP = \{\tt{FILTER}, \tt{GROUP}, \tt{BACK}\}$). Each operation $o \in OP$ is defined as follows,
\begin{itemize}
    \item $\tt{FILTER(attr,op,term)}$ is used to select data tuples that match a criteria. It takes an attribute, a comparison operator (e.g. equals, contains) and a numerical/textual term, and results in a new display representing the corresponding filtered subset.
    \item $\tt{GROUP(g\_attr, agg\_func, agg\_attr)}$ is used to group and aggregate the data. It takes a single attribute to be grouped by, an aggregation function and another attribute to employ the aggregation function on.
    \item $\tt{BACK()}$ allows a backtrack to the previous display in order to enable an alternative exploration path.
\end{itemize}
This query is encoded into a six dimensional vector. The encoded query, $q_t$ can be represented as $q_t = \{o, attr, op, term, g\_attr, agg\_func\}$.
We design the state space of our model such that it can capture the history (queries and their results) of EDA session along with the intent identified, $I_t$ from the BTM model and the cost incurred $C_t$ by the agent. Figure~\ref{} illustrates different components of our state space. To preserve the broader context in which the results are generated, we concatenate last two query ($q_{t-1}, q_{t-2}$) and display vectors ($v_{t-1}, v_{t-2}$) to the current query and display vector. The state space can be thus represented as (\ref{eq:state_space}).


\subsubsection{Action Space}
At each time-step t, for the query, $q_t$, our agent outputs an action $a \in A$. Each of this action corresponds to a downsampled dataset $a_i$. Each action has a corresponding cost of running $c_i$, represented by the fraction of rows processed by the agent while taking that action. 
\sm{Notational explanation. Details on how it goes back to the EDA sequence.}

\subsubsection{Neural Network Architectures}
The A2C model has two components - an actor (policy based) and a critic (value based). The actor takes in the input and outputs the best action whereas, the critic evaluates the action by computing the value function. We model the actor (policy network) as a function approximator using a fully connected neural network. The critic (value network) is also separately modelled by a fully connected neural network. The critic receives the state and the action from the actor as its input and it outputs the Advantage value. 
The training for both the networks happen separately, using gradient descent to find the global maximum and update their weights. As time passes, the policy network tries to learn the best policy and the value network gets better at evaluating those actions.

\end{comment}
\begin{comment}
\sm{A2C -- what is A2C? How many networks, how does updates work.
For us what is the high-level architecture.}

\sm{Refer to the Figure to explain each component of state Space}
\end{comment}

% \subsection{Training Strategy}
% \sm{Give the detailed algorithm for RL using A2C showing updates.}

% \tong{This section is called model overview, but no model details (e.g., parameterization) are introduced. Please see an example in Section 3.1 of https://arxiv.org/pdf/1805.00145.pdf.}
% \sg{Explain our setup for the generation of the EDA sequences, so that we can further in the paper use it to motivate what we are trying to do. This will include the limitation that we don't have humans to generate the sequences and hence we are trying to model humans via this simulator from the ATENA paper.\\
% After this we can motivate our problem by showing some sequences from the simulator on a fixed sample and the original table and show how they lead to different insights.\\
% \colorbox{yellow}{Experiment section 1:} - Need sample sequences to motivate our introduction, we already have 1, we need 1 more. What we are trying to do here is we will show that there is a divergence in the sequence. We fix the first 3 sequences and then let the simulator run its course on the rest. We have to show the results qualitatively as shown in the demo.
% Give an overview of the solution approach.}

\noindent\textbf{Intent-divergence vs. intent-shift}
This paper addresses the problem of \textit{intent-divergence} where the users' have a set of consistent implicit objectives/intents during the EDA. 
While \name{} deliberately introduces a bias towards a set of intents inferred from historical analysis sequences, this is justified as real expert users often follow some standard exploration strategies and look for similar types of insights even on new datasets as observed by multiple prior works~\cite{atena2020sigmod, jain2016sqlshare, milo2018next} and do not perform random explorations.
However, if such implicit objectives suddenly change causing an \textit{intent-shift}, then \name{}'s current design would not be able to detect that. 
In this paper, it is assumed that there is limited \textit{intent-shift} (\emph{i.e.}, change of user's core intent characteristics). 
%While in this paper we are the first to address how to prevent intent-divergence (approximation errors misleads users to a different intent), 
This assumption is widely adopted in other interactive applications \cite{liu2018dialogue,guo2018dialog,tan2019drill}, although it is interesting to address the intent-shift \cite{xie2021tiage,arnold2019sector}. As the first work on approximate exploratory data analysis in the interactive setting, we leave addressing \textit{intent-shift} as future work.

%\vspace{-0.5em}
\subsection{Implementation \& Deployment}

\begin{algorithm}[th]
        \begin{algorithmic}[1]
        \label{alg:alg1}
        \caption{\name{}'s Training Algorithm}
        % \textbf{Input} Dataset $D$
        \STATE Train the ATENA~\S{5.2} on given dataset $D$
        \STATE Generate sequences $O$ using ATENA
        \STATE Train a $BTM$ model $\phi$ with $K$ topics on $O$
        \STATE Create sample set $A=\{a_1,a_2,...,a_i\}_{i=1}^n$
        \STATE Initialise A2C model
        \FOR{$i = 1,...,episodes$}
            \STATE Initialize the state space $s_t$ as in (\ref{eq:state_space}) with zeroes
            \FOR{$t = 1,...,l$}                    
                    \STATE Generate a query $q_t$ from the simulator
                    \STATE Take action $a_t$ for the given $s_t$
                    \STATE Update $s_t$ with latest query $q_t$, display vectors $v_t$, latency $C_t$ and current intent as $\phi(Q_t)$
            \ENDFOR
            \STATE Calculate Individual rewards using~\S{4.3}.
            \STATE Calculate the total reward $R_t$
            \STATE Calculate $J_R(\pi) = \mathbb{E}_{a_t \sim \pi_\theta}[R_t]$
            \STATE Compute loss functions $\mathbb{L}(\theta)$ for policy network
            \STATE Compute $\mathbb{L}(\omega)$ for value network
            \STATE Update A2C model
        \ENDFOR
        \end{algorithmic}
    \end{algorithm}

%\vspace{-0.2em}
Algorithm 1 explains the steps we use for training our RL Agent. It takes in the dataset $D$, processes it by creating samples as shown in line 4 and training ATENA simulator on the same as shown in line 1. Then on the generated sequences, a BTM model is trained in line 3. Once this is done, we then start the training of our agent as demonstrated in lines 5-18 of Algorithm 1. We follow episodic training where we calculate rewards at the end of each episode as shown in line 13. Once we calculate the total reward, we compute the loss functions and update both actor (policy) and critic (value) networks as explained in \S{\ref{apps:model_implementation}}. Appendix \textbf{\ref{app:discussion}} discusses deployment aspects.


%This limitation is widely applicable to other intent-based systems as well~\cite{} we leave this as a future work.
%\sm{Tong - you have any references?}\tong{Look}
%\sm{“The authors do not sufficiently discuss the downsides of an intent-aware sampling approach. For example, the proposed strategy introduces a bias towards, unlike its main competitors that sample uniformly at random. Consequently, it remains unclear whether the proposed method can estimate a truly exploratory user intent, as it will be more random (ln. 118). The introduced bias appears to be more problematic than the accumulated approximation error.”}

%\sm{“The authors introduce a new concept of intent-divergence early on (ln 164), but however, never explain it in any detail. It remains unclear why the proposition prevents an intent-divergence of the user. If the analyst’s intent actually changes, preventing an intent divergence on the model side, seems to hinder their workflow.The manuscript does not include a discussion of the relevant problem of “intent-shift”.”}


\begin{comment}
Training of this model requires us to have human generated sequences of queries on a given dataset and these sequences not readily available. So for the purpose of this paper, we are generating these query sequences using a simulator proposed by \cite{atena2020sigmod}. This simulator mimics a human expert trying to do EDA on a given data and generates sequences, $Q$ of fixed length $l$.
For each query, the EDA operations considered are \emph{Filter}, \tong{we may use \tt{Filter}} \emph{Group} and \emph{Back} where,
\begin{itemize}
    \item \emph{FILTER(col, op, term)} is used to select data tuples that match a criteria. It takes an column, a comparison operator (e.g. equals, contains) and a numerical/textual term, and results in a new display, the corresponding filtered subset.
    \item \emph{GROUP(g\_attr, agg\_func, agg\_attr)} is used to group and aggregate the data. It takes a single attribute to be grouped by, an aggregation function and another attribute to employ the aggregation function on.
    \item \emph{BACK()} allows a backtrack to the previous display in order to enable an alternative exploration path.
\end{itemize}
Each query $q_i$ of the sequence $Q$ is represented as a six dimensional vector
\begin{equation}
\begin{multlined}
    q_i = [Operation Type, Column ID, Filter Operator,\\
     Filter Term, Agg Column ID, Agg Function] 
\end{multlined}
\end{equation}
\tong{I would suggest to define notation for each of them in the above.}
Here $OperationType$ takes on values $\{0,1,2\}$, where 0 represents back, 1 represents filter and 2 represents group type query.

The set of EDA sequences, $QS = \{Q_1, Q_2, .., Q_j \}_{j=1}^m$ \tong{lower case for consistency} are generated on the original dataset $D_0$ and are considered to be the original EDA sequences which mimic a data analyst. These sequences are then fed into the intent model, which uses topic modelling to cluster these sequences into different topics. \tong{this part needs to be presented more formally. I would suggest to follow some existing papers, instead of starting from scratch.} Each of these topics corresponds to different intent. 
We then train our RL agent on these generated sequences and the identified intents using a novel reward function which takes into account the intents and the given latency, and chooses an action from the given downsampling strategies. Figure \ref{fig:model_overview} explains how our model will be used by an analyst who wants to do EDA. \sg{Add a figure for the model architecture}
\end{comment}

\begin{comment}
\sg{Include a figure for the reference of training.}

\begin{figure*}[t]
\centering
\begin{minipage}{0.5\textwidth}
\includegraphics[width=1.0\linewidth]{figures/solution_outline.png}
     \caption{Model Overview. \tong{we may use pdf or ps format for the figure (instead of png), to get a higher resolution figure? Besides, this might be 'system overview'. Intent-ware is a key point, but is downplayed in the figure for now. Where is the topic model?}}
     \label{fig:model_overview}
\end{minipage}%
\hfill
\begin{minipage}{0.5\textwidth}
\centering
%\includegraphics[width=\linewidth]{figures/framework.png}
\caption{}
\label{fig:framework}
\end{minipage}
\end{figure*}

\end{comment}

\begin{comment}

%\subsection{Intent-Aware Reinforcement Learning}
\subsection{Reward Design}
\tong{}
Reinforcement learning aims to learn an optimal policy for an agent interacting with an unknown and often a complex environment. A policy is defined by a conditional distribution $\pi(a|s)$, which denotes the probability of choosing an action $a \epsilon A$ when agent is at state $s \in S$ If the agent chooses an action $a \in A$ at state $s \epsilon S$, it receives an immediate reward $r(s,a)$. To solve our problem, we introduce three different kinds of rewards, latency reward, intent reward, and termination reward to refine the policy from different perspectives.

\subsubsection{Latency Reward}
In an interactive exploratory data session, we want users to get instant results. Hence, we want out agent to choose the samples from the sample space that has least number of rows. At a state $s_t$, when an agent takes an action $a_t$, the total number of rows read by the agent while taking that action is termed as $|{T_s}|$ and the total number of rows in the original dataset as $|T|$. Then the immediate reward after taking an action $a_t$ is,
\begin{equation}
    r(s_t, a_t) = 1 - \frac{|{T_s}|}{|T|}
\end{equation}
\begin{equation}
    R_{latency} = \sum_{t=1}^{l} r(s_t, a_t)
\end{equation}
where $l$ is the length of the sequence. This reward ensures that the agent does not always select the sample with the most amount of data, as that data will lead to best approximation of data.
\subsubsection{Intent Reward}
While choosing the smallest samples using the latency reward will lead to reduced waiting time for the user, but it will often lead to diverging of intent after some point. If we only use latency reward, agent will often always choose the smallest sample out of the all available. We not only want the agent to choose samples, that will minimize time, but we also want it to preserve the intent of the user. To do so, we introduce two components in the intent reward. We denote the distance function for distance between two query sequences to be $Distance(Q_1, Q_2)$. We define the reward function to be,
\begin{equation}
\label{eq_eda_sim}
    R_{dis} = 1 - min(\{Distance(Q_t, O_i)\}_{i=1}^{m})
\end{equation}
Here $m$ is the length of the sequences on the original table. For our purpose, we use the distance metric, EDA-Sim defined in \cite{edasim2018kdd}. EDA-Sim compares the order and the content of the results from the query sequence and gives a similarity score. We use this metric to find the closest sequence from $O$, so that we can match our trajectory with one of the sequences from the original sequences. For the rest of this paper, this identified closest sequence can be referred to as $Q_{ground}$

However, only this reward is not sufficient for preserving the intent. There are cases where query sequences with different intents may have a low EDA-Sim value. Hence to better match the intent, we define a reward, $R_{topic}$, which is,
\begin{equation}
    R_{topic} = 1 - D(\phi(Q_t), \phi(Q_{ground}))
\end{equation}
For our purpose, the $D$ function used is Euclidean Distance Measure. The total reward can thus be defined as,
\begin{equation}
    R_{intent} = R_{dis} + \delta*R_{topic}
\end{equation}
These rewards ensures that the topic distribution matches with that of the closest sequence.
\subsubsection{Termination Reward}
In an EDA sequence, last few queries are the most important. These generally contain the results of the whole EDA session and are generally the ones that has insights. Even after, matching intent of the user with the closest sequence, it might not be enough to preserve these insights. Therefore, the third component of our reward tries to make sure that these insights are preserved. The first part of this reward looks at the query sequences and not their results. This is a binary reward. For a given sequence $QS_t$, we define this reward as follows:
\begin{equation}
    R_{match} = \{0, 1\}
\end{equation}
This reward is 1, if the last $k$ queries of the sequence generated by the agent match with any of the sequence in the original sequences, that belong to the same topic. The other component of the reward looks at the results of the queries instead of the query definition. It is important for us to have a reward defined on the results of the last $k$ queries because different queries might lead to similar insight. For a given generated sequence, $Q_t$ and the closest sequence identified, $Q_{ground}$ in \ref{eq_eda_sim}, we define this reward as,
\begin{equation}
    R_{recall} = top\_k\_recall(Q_t, Q_{ground})
\end{equation}
The final termination reward can be written as,
\begin{equation}
    R_{term} = R_{match} + \zeta*R_{recall}
\end{equation}
With all the three reward functions, the reward function $J_R(\pi)$, can be written as,
\begin{equation}
    J_R(\pi) = \mathbb{E}_{a_t \sim \pi_\theta}[R_t]
\end{equation}
where the Reward function $R_t$ is,
\begin{equation}
    R_t = R_{latency} + \alpha*R_{intent} + \beta*R_{term}
\end{equation}. 
%To train the RL model with the given reward function we use the policy-gradient bases Advantage Actor Critic (A2C) with neural network as a function approximator for the policy and value functions. 
%For the policy network and the value network, we have different neural networks, each with 2 hidden layers of latent dimension 64.
After learning the optimal poilcy $\pi_\theta(a|s_t)$, the action at each step is chosen via,
\begin{equation}
    a_t = argmax_{a \in A}\pi_\theta(a | s_t)
\end{equation}

\end{comment}
% Explain the state space and action space here. 
% Have a pseudocode here for the RL part.\\
% \tong{Instead of only RL, we may create a pseudocode for the overall system. Please see an example, Algorithm 1, in https://aclanthology.org/N18-1187.pdf (NAACL)}
% \tong{I think it is fine even if we do not have modifications to the RL algorithm. Just clearly define everything. Here is an example in Section 4.2.3 of https://arxiv.org/pdf/2007.00194.pdf (KDD). Another example is Section 3.3. in https://arxiv.org/pdf/1811.10092.pdf (CVPR best student paper). Another example is Section 3.7 in https://aclanthology.org/N18-1187.pdf (NAACL). Another example is Section 3.2.2 in https://arxiv.org/pdf/1805.00145.pdf (NeurIPS)}
% \colorbox{yellow}{Experiment section 3:} - RL training changes:
% \begin{itemize}
%     \item Increase the state space of the experiment by including latency in the state space.
%     \item Reward Design - Have all the components discussed in here . 
%     \item On a batch of sequences, have a penalty term to make the agent generate sequences closer to the original distribution
%     \item Vary episode lengths for 10,15,20 and compute error metrics like order and recall error.
% \sg{
% \begin{enumerate}
%   \item EDA-SIM Select Min Distance from heldout set 
%     \item Latency Reward 
%     \item Prob. Distribution over topics vector euclidean Reward for each sequence
%     \item Cumulative topic distribution euclidean distance reward 
%     \item Last k rows query matching 
%     \item Top $p$ precision for the last k query results (EDA-SIM) 
%     \item Top $p$ recall for the last k query results (EDA-SIM) 
%     \item TODO: Confidence interval (Show time savings)
% \end{enumerate}
   
% \begin{equation}
%     Reward = \alpha*(1+(6+7)/2) + \beta*4 + \gamma*2 + \delta*3 + \zeta*5    
% \end{equation}

% }
% \end{itemize}

% Have a table about all the samples and the size of the original dataset for each dataset

% TODO: Data: 500k rows, 350k training simu , 150k testing simu (TODO)
% TODO: Query Column Sets
% TODO: All rewards between 0 to 1

% Ablation Study Experiments:
% Categorize rewards into three categories

% 2 sets of experiments, full trained, 4k trained

% Step 1: Generate 5k sequence on one data from simulator
% Step2: Take out 4k/5k to train our RL, heldout is 1k
% Step3: Training of the Agent
% Step4: Generate 1k sequences from our RL agent (Simulator + Agent)
% Step5: Evaluation

% Metric 1: Compare step4 with step2 heldout (intent accuracy)
% Output: Baseline (Fixed, Baselines, Our agent), Top1, Top2, Top3
% Metric 2: Compare latencies (Percentage)
% Output: Both seconds + Actual number of rows (summation)
% Metric 3: Seq in step 4, run those on original data and then calculate the precision and recall for the groups in the result

% For the set of sequences of ~5400 that we have, we take out 30 sequences For each topic, choose 3 randomly.\\


% 1. For each seq in heldout set we have prob. dist. over topics
% 2. 

% T1 ...  T10

% Heldout (A):

% Session 1: [0.1 , 0.05,  ....   0.25]. => arg([]) > t => [T3,  ]
% Session 2: [0.1 , 0.05,  ....   0.25] =>  [T5,  ]
% .
% .
% 1000 --> 

% RL-intervened (B):
% Session 1: [0.1 , 0.05,  ....   0.25]. arg([]) > t =>  [T3,  ]
% Session 2: [0.1 , 0.05,  ....   0.25] =>  [T3,  ]
% Session 3: [0.1 , 0.05,  ....   0.25]

% .
% .
% .
% 1000

% Metric:


% count(Tx) in A
% ----------    
% 1000. filter (>t) -> ??

% vs

% count(Tx) in B
% ----------     
% 1000